\clearpage
\section{Lift Waiting Location}
The code solving this, with in depth documentation, is attached as q1.py, the result it produces is floor 4.
\subsection{Bonus Calculation}
To find the optimal resting floor we'll start by analyzing the paths it takes, focusing on a single person. Through all the assumptions we can simplify the model to:
\begin{enumerate}
	\item Person enters on floor 0.
	\item Person go to floor X.
	\item Person then goes to 3 random floors.
	\item From the final floor he reaches, the person goes down to 0 and leaves.
\end{enumerate}
in each movement step the elevator performs 
\begin{enumerate}
	\item Elevator goes to origin floor from its resting point.
	\item Elevator goes to destination floor from the origin floor.
	\item Elevator goes to the resting point from the destination floor.
\end{enumerate}
Each ride's cost calculation is 
\begin{equation}
	C=\abs{rest-origin}+\abs{dest-origin}+\abs{rest-dest}
\end{equation}
we can notice that the middle element doesn't depend on the rest point so we can simply ignore it. Therefore the goal we want to minimize is 
\begin{equation}
	C=E\ins{\abs{rest-origin}}+E\ins{\abs{rest-dest}}
\end{equation}
but from the symmetry of a sufficiently large random selection of people and trips we can simplify this to 
\begin{equation}
	C=2E\ins{\abs{rest-X}}
\end{equation}
where X is a floor. To this we add the trips as the mapping 
\begin{equation}
	\{0,f_1,f_2,f_3,f_4\} \to \{f_1,f_2,f_3,f_4,0\}
\end{equation}
we notice that the floor \(0\) appears twice and because the person travels from/to there in two separate trips, so when working with the floor numbers, we'll need to weight trips to/from the floor \(0\) at double the weight of the others. We can now calculate
\begin{equation}
	E\ins{\abs{rest-X}}=\frac{rest}{5}+\frac{4}{5}\cdot\frac{1}{10}\sum_{X=1}^{10}\abs{rest-X}
\end{equation}
we take the derivative to find the optimum
\begin{equation}
	\dv{E}{rest}=\frac{1}{5}+\frac{4}{50}\sum_{X=1}^{10}\text{sgn}\ins{rest-X}
\end{equation}
and 
\begin{equation}
	\frac{4}{50}\sum_{X=1}^{10}\text{sgn}\ins{rest-X} + \frac{1}{5}=0
\end{equation}
Which is a pain to solve so we turn to tricks. If we assume there are \(k\) floors below the resting point, this means the expression inside the sum will be \(+1\) \(k\) times, and \(-1\) \(10-k-1\) times. Therefore 
\begin{equation}
	D=\frac{4}{50}\sum_{X=1}^{10}\text{sgn}\ins{rest-X} + \frac{1}{5} = \frac{1}{5}+\frac{4}{50}\ins{k+k-9}=\frac{10-36+8k}{50}=\frac{8k-26}{50}
\end{equation}
plugging in the different possible values for k
\begin{align}
	D(0)=-\frac{26}{50}&&D(1)=-\frac{18}{50}&&D(2)=-\frac{10}{50}&&D(3)=-\frac{2}{50} && D(4)=\frac{6}{50}
\end{align}
so our two suspect floors are \(3\) and \(4\). We look at the values at the floors:
\begin{align}
	E\ins{\abs{3-X}}&=\frac{3}{5}+\frac{4}{50}\sum_{X=1}^{10}\abs{3-X}=\text{calculator}=\frac{77}{25}\\
	E\ins{\abs{4-X}}&=\frac{4}{5}+\frac{4}{50}\sum_{X=1}^{10}\abs{4-X}=\text{calculator}=\frac{74}{25}
\end{align}
so floor 4 gives a lower expected value, meaning it is our optimal floor.